\documentclass[12pt]{article}
\usepackage[T1]{fontenc}
\usepackage[utf8]{inputenc}
\usepackage{lmodern}
\usepackage[margin=1in]{geometry}
\usepackage{amsmath,amssymb}
\usepackage{setspace,microtype}
\usepackage{booktabs,longtable,array,tabularx}
\usepackage{graphicx}
\usepackage{url,xurl}
\usepackage[round,authoryear]{natbib}
\usepackage{caption}
\usepackage[colorlinks=true,linkcolor=blue,citecolor=blue,urlcolor=blue]{hyperref}
\usepackage{bookmark}
\newcolumntype{P}[1]{>{\raggedright\arraybackslash}p{#1}}
\setlength{\emergencystretch}{3em}
\setlength{\parskip}{0.3em}
\setcounter{tocdepth}{2}
\captionsetup{font=small,labelfont=bf}
\urlstyle{same}
\let\harvardurl\url
\hypersetup{
 pdftitle={Switchback Experiments: Design, Inference, and the Credibility of Carryover Assumptions},
 pdfauthor={},
 pdfkeywords={temporal interference, experimental design, randomization inference, causal estimands, carryover},
 pdfcreator={LaTeX}}
\title{\textbf{Switchback Experiments: Design, Inference, and the Credibility of Carryover Assumptions}}
\author{}
\date{October 2026}
\bibliographystyle{agsm}

\begin{document}
\singlespacing
\maketitle
\begin{abstract}
Switchback experiments repeatedly randomize treatment over time, often within an aggregate marketplace or operational unit. They can reduce interference between simultaneously treated and untreated users, but create difficult questions about how historical assignments affect current outcomes. This focused review synthesizes 34 supplied papers and related methodological studies, emphasizing their stated carryover assumptions, the quantities those assumptions require researchers to know or choose, and the causal targets their procedures actually address. We distinguish exact finite memory, geometric mixing, declining history sensitivity, observed-state models, linear impulse responses, and simulator-based design. We also separate assignment optimization, estimation, confidence intervals, hypothesis tests, diagnostics, precision studies, and rollout decisions. The comparison identifies a recurring inferential gap: a local-exposure estimator can be asymptotically normal while retaining bias for a permanent-policy effect. The memory scale is sometimes supplied as a design input, sometimes constrained only through an existence assumption, and sometimes replaced by stronger structural information. Credible applications should justify these commitments separately, report sensitivity to their calibration, and distinguish absence of detected carryover from evidence that a chosen horizon is sufficient.
\end{abstract}
\noindent\textit{Keywords:} temporal interference; experimental design; causal inference; dynamic treatment effects; carryover; randomization inference.
\clearpage
\tableofcontents
\clearpage
\setstretch{1.8} % DON'T change the spacing!

\section{Introduction and scope}\label{sec:introduction}

Switchback experiments assign a treatment repeatedly to the same system over time. They are especially useful when individual randomization is infeasible or when users compete for shared resources, so that treating one user changes another's outcome. Aggregating an interacting population into a marketplace, region, or item cluster can make the treatment contrast easier to interpret. However, yesterday's treatment can affect today's inventory, supply, behavior, or system state. Switching the assignment does not necessarily switch the outcome-generating process immediately.

The literature consequently addresses several different problems under the same experimental label. Some papers choose assignment sequences to reduce worst-case estimation error. Others estimate dynamic or long-run causal effects, construct uncertainty statements, diagnose a misspecified carryover horizon, compare power, or make deployment decisions. Their temporal assumptions also differ. An exact last affected lag, a mixing time, a stable autoregressive coefficient, and an observed sufficient state express different restrictions. Calling each of them a carryover parameter obscures what must be supplied before an experiment and what the data can actually learn.

This review is organized around those distinctions. It examines the 34 PDFs supplied in the project's literature collection, including closely related time-series, panel, crossover, and dynamic-policy methods. It is a focused review of that collection, not an exhaustive systematic search or quantitative meta-analysis. Following the review protocol, we read introductions, related-work sections, and discussions or conclusions, then made narrow lookups of assumptions, estimand definitions, and qualifications needed to reproduce the carryover restrictions accurately. We do not independently verify proofs or rank all methods by a common empirical benchmark.

The unit of comparison is the reviewed manuscript version. Publication years in citations do not always coincide with the local PDF's revision date. In particular, \citet{huwager2025} is the December 2025 fourth version; \citet{nipanel2024} is the 2024 Harvard working paper rather than its revised 2026 successor. The local copies underlying \citet{basse2023} and \citet{liangrecht2025} are respectively 2020 and 2024 preprints of later journal publications. Several files catalogued as 2025 contain 2026 revisions, recorded in the bibliography. The local \citet{ni2022} submission is anonymized; authorship follows the \href{https://sites.google.com/view/caudyn2022/accepted-papers}{official workshop programme}. These distinctions matter when later versions change estimands or inferential guarantees.

\section{A framework for comparing claims}\label{sec:framework}

\subsection{An experimental design does not determine a unique estimand}

For exposition, let $W_t\in\{0,1\}$ denote the treatment assigned at time $t$, and let $Y_t(w_{1:t})$ be a nonanticipating potential outcome. This is normalized review notation, not an assertion that every source imposes nonanticipation in the same way. A finite-horizon global treatment effect is
\begin{equation}\label{eq:global-effect}
 \tau_T=\frac{1}{T}\sum_{t=1}^{T}
 \left\{Y_t(\mathbf{1}_{1:t})-Y_t(\mathbf{0}_{1:t})\right\},
\end{equation}
with a specified common initial condition or pre-experiment history. This compares policies applied throughout the observed horizon. In a panel, one can additionally average across units. Spatial interference requires indexing outcomes by the relevant treatment matrix, not merely the unit's own treatment.

A recent-run contrast instead changes the last $q$ assignments while keeping, conditioning on, or averaging over an earlier history. General lagged incremental effects can change a single assignment and average over subsequent assignments. Neither construction automatically equals Equation~\eqref{eq:global-effect}. Equality needs a suitable exposure restriction or an approximation argument. A stationary reward contrast concerns policy-specific equilibrium distributions; a discounted value contrast weights future outcomes from a reference initial state; a local policy derivative measures an infinitesimal change. These are useful targets with different interpretations, not interchangeable estimates of one universal effect.

A simple review example makes the distinction concrete. Suppose
\[
 Y_t=\mu_t+\beta_0W_t+\beta_1W_{t-1}+\epsilon_t,
\]
with independent fair assignments and mean-zero noise independent of them. Comparing current treated and control observations targets $\beta_0$, because the previous assignment is balanced in expectation. Sustained treatment versus sustained control, after the initial boundary, has effect $\beta_0+\beta_1$. A precise estimate of the first contrast can be misleading for the second even though treatment is randomized. This example is illustrative rather than a result attributed to a particular source.

\subsection{Three distinct knowledge requirements}

A method may require a scientific restriction without requiring its numerical parameter as an estimator input. For example, a central limit theorem can assume some finite mixing constant exists, while computing an average needs no such value. A recommended burn-in that explicitly multiplies that constant is a different matter. Conversely, a chosen truncation length is an analyst's tuning parameter, not proof that effects vanish beyond it. We therefore distinguish maintained existence assumptions, operational inputs, and estimated nuisance quantities throughout the review.

The distinction extends to spatial structure. Aggregating users can reduce within-market interference while leaving dependence between geographic or item clusters. A paper using multiple regions may still assume no spillover between them. Only methods that specify and handle such exposure warrant a spatial-interference classification. Similarly, autocorrelation in untreated outcomes is not itself treatment-induced carryover: persistent weather shocks and treatment-induced resource changes can produce superficially similar time series with different causal implications.

Table~\ref{tab:comparison} records the main contribution and endpoint of each supplied study. D denotes design, A analysis, and D+A both. Endpoint codes are E for causal estimation, C for confidence intervals, H for hypothesis testing, V for variance or estimation-error analysis, P for power, and R for a deployment decision. They summarize the contribution emphasized in the reviewed version; a design paper may use an estimator without establishing a new general inference procedure. The detailed sections supply the restrictions behind each abbreviated entry.

\begin{singlespace}
\begingroup
\footnotesize
\setlength{\tabcolsep}{3pt}
\renewcommand{\arraystretch}{1.18}
\begin{longtable}{P{3.2cm}P{0.9cm}P{1.8cm}P{4.5cm}P{1.3cm}P{3.3cm}}
\caption{Contribution, interference scope, and inferential endpoint in the supplied collection.}\label{tab:comparison}\\
\toprule
Study & Work & Scope & Carryover restriction or required information & Output & Target or qualification\\
\midrule
\endfirsthead
\multicolumn{6}{l}{\textit{Table \thetable\ continued}}\\
\toprule
Study & Work & Scope & Carryover restriction or required information & Output & Target or qualification\\
\midrule
\endhead
\midrule
\multicolumn{6}{r}{\textit{Continued on next page}}\\
\endfoot
\bottomrule
\endlastfoot
\citet{bojinov2023} & D+A & Temporal & management science. Exact finite order $m$; specify $p$; pilot comparisons proposed & E,C,H & Persistent contrast; underestimating memory biases it\\
\citet{liuzhong2026} & A & Temporal & same assuption as \cite{bojinov2023} & E,H & Partially sharp and weak nulls; lower bound on persistence\\
\citet{ni2022} & D+A & Temporal & 1-carryover. preliminary version of \cite{bojinov2023} presented at NeeurIPS worksop; period length chosen to cover persistence & E,C,H,V & Balanced single-unit exposure\\
\citet{masoero2023} & D+A & None in baseline & Amazon. No cross-unit or temporal interference in central analysis & E,V & Balanced unit--time ATE\\
\citet{missault2025} & D+A & Temporal + item & Amazon. $m$-carryovere; no cross-item spillover & E,V,P & Regular balanced multi-unit exposure\\
\citet{zeng2026} & D+A & Temporal & Airbnb. Zero or exact first-order carryover; no spillovers & E,C,H & Exact H without carryover; asymptotic C with carryover\\
\citet{chen2026} & D+A & Temporal & management science. $m$-carryover, sufficient surrogate that captures carryover, stable conditional densities, divergence bound & E,C,H & Oracle and plug-in guarantees differ\\
\citet{xiong2024} & D+A & Temporal; simultaneous policies & Kernel formulation that generalizes finite carryover and decaying model; prior cumulative-effect curves & E,V & Global effect holding other policies fixed\\
\citet{ni2025model} & D+A & Temporal & Decreasing history envelope; chosen clock, balance and coupling & E,V,C & CLT centered at estimator expectation\\
\citet{ni2025decision} & A & Temporal extension & Outcome smoothness and diminishing perturbation effects & R & Design given; minimax rollout regret\\
\citet{carryover2025} & A & Temporal & Additive lag effects for diagnostic theory; candidate lag chosen & H & Reject insufficient lag; nonrejection is not certification\\
\citet{huwager2025} & D+A & Temporal & Uniform geometric contraction; calibrated burn-in and blocks & E,C,V & Filtered effect versus corrected GATE\\
\citet{jia2025} & D+A & Temporal + spatial & Geometric contraction, graph, exposure restrictions & E,V & Global policy effect; graph-dependent rates\\
\citet{hu2026} & D+A & Temporal & Exact finite $m$ supplied to duration design & E,V & Spectral MSE and random run durations\\
\citet{glynn2020} & D+A & Temporal & Observed finite state, stationary irreducible chains, exploration & E,V & Difference of stationary rewards\\
\citet{blanchet2026} & A & Temporal & Geometric convergence plus monotonicity restrictions & E,V & Transient versus stationary effect; conservative inference conditions\\
\citet{wen2025} & D+A & Temporal & Observed Markov state, repeated daily trajectories, estimator-specific stability & E,V & Cadence depends on carryover and residual dependence\\
\citet{farias2022} & A & Temporal & Observed state, mixed-policy mixing, small transition perturbation & E,V & Stationary effect with approximation bias\\
\citet{shi2023} & A & Temporal & Sufficient stationary Markov state and geometric ergodicity & E,C,H & Discounted value contrast; sequential testing\\
\citet{li2026} & D+A & Temporal & Structural finite-capacity queue; growing interval lengths & E,C & Local stationary throughput derivative\\
\citet{sun2024} & D & Temporal & Stable controlled ARMA; orders and coefficients estimated & E,V & Small-signal design optimality\\
\citet{wu2026} & D & Temporal & Full observed history and counterfactual simulation environment & V & Learned allocation for supplied estimator\\
\citet{johari2025} & A & Temporal & Uniform contraction; truncation, perturbation and mean-stability restrictions & E,C,V & GATE approximation; expectation-centered CLT\\
\citet{bojinovshephard2019} & A & Temporal & General history; nonanticipating assignment; positivity & E,C,H & Lagged contrasts need no finite cutoff\\
\citet{bojinovpanel2021} & A & Temporal & General own-unit history; no cross-unit interference & E,C,H & Dynamic panel causal effects\\
\citet{linding2025} & A & Temporal & Independent assignment; coupling decay or growing finite memory & E,C,H & Lag effects; conservative HAC inference\\
\citet{liangrecht2025} & A & Temporal & Linear invariant impulse response; tail and concentration conditions & E,C & Response functionals, not arbitrary histories\\
\citet{guoliang2026} & D & Temporal & Fixed finite impulse-response order $K$; model-based design & E,V & Target-dependent Markov assignments\\
\citet{jiangding2025} & A & Temporal & Sequence support; optional finite memory and effect restrictions & E,C,H & Crossover identification; restriction tests\\
\citet{basse2023} & D & Temporal & Main habituation model; finite pulse decay only for recycling & E,V & Pulse/wedge targets; not all-history cutoff\\
\citet{nipanel2024} & D+A & Temporal + spatial & One-period memory and known spatial neighborhood & E,C & Local 2024 panel exposure model\\
\citet{ferrari2026} & A & None in generator & Contemporaneous effect; calibrated serial correlation & V,P & Applied cadence/precision comparison\\
\citet{pankratev2026power} & A & None in baseline & Contemporaneous treatment and multilevel shocks & V,P & Cluster-size imbalance and power\\
\citet{pankratev2026variance} & A & Temporal + spatial in simulation & Specified three-lag and spatial contamination generators & V,P & Adjustment comparison; no general interference remedy\\
\end{longtable}
\endgroup
\end{singlespace}

\section{Carryover assumptions and methods for switchback experiments}\label{sec:carryover}
The literature distinguishes two broad carryover assumptions. Under finite memory, assignments before a fixed horizon have no effect on the current outcome. Under decaying memory, earlier assignments can remain influential, but their effects diminish with time.
\subsection{Design under finite and approximate carryover}\label{sec:finite}

\subsubsection{$m$-carryover effects}

\citet{bojinov2023} develop a design-based framework for switchback experiments with a known finite carryover order. Their Assumption~2, ``$m$-Carryover Effects,'' begins ``There exists a fixed and given $m$'' \citep[Assumption~2, p.~3762; PDF p.~4]{bojinov2023}. It requires
\begin{equation}\label{eq:finite-memory-original}
Y_t(w'_{1:t-m-1},w_{t-m:T})
=Y_t(w''_{1:t-m-1},w_{t-m:T}),
\qquad t=m+1,\ldots,T.
\end{equation}
Combined with nonanticipation, this means that the current outcome depends only on the current treatment assignment and the preceding $m$ assignments, $w_{t-m:t}$. Earlier assignments have no effect; the assumption does not restrict the magnitude or form of effects within this window.

The paper distinguishes a sharp null from a weak null. Under the sharp null, each analyzed outcome is unchanged by the assignment path, permitting Fisher randomization inference. The weak null specifies zero average effect and allows positive and negative period-specific effects to cancel. With $p=m$, the average effect is
\begin{equation}\label{eq:finite-memory-estimand}
\tau_p(Y)=\frac{1}{T-p}\sum_{t=p+1}^{T}
\left[Y_t(\mathbf{1}_{p+1})-Y_t(\mathbf{0}_{p+1})\right],
\end{equation}
where $\mathbf{1}_{p+1}$ and $\mathbf{0}_{p+1}$ denote consecutive treatment and control assignments over $p+1$ periods \citep[Eq.~(1), p.~3762]{bojinov2023}. The weak null is $H_0:\tau_p(Y)=0$.

The Horvitz--Thompson estimator weights observed constant-treatment runs by their probabilities under the experimental design:
\begin{equation}\label{eq:finite-memory-estimator}
\widehat{\tau}_p
=\frac{1}{T-p}\sum_{t=p+1}^{T}\left[
\frac{Y_t^{\mathrm{obs}}\,\mathbf{1}\{w_{t-p:t}=\mathbf{1}_{p+1}\}}
     {\Pr(W_{t-p:t}=\mathbf{1}_{p+1})}
-\frac{Y_t^{\mathrm{obs}}\,\mathbf{1}\{w_{t-p:t}=\mathbf{0}_{p+1}\}}
      {\Pr(W_{t-p:t}=\mathbf{0}_{p+1})}
\right].
\end{equation}
Here $\mathbf{1}\{\cdot\}$ is an indicator, and the exposure probabilities are known from the randomization scheme and must be positive. This reproduces Equation~(4) with design arguments suppressed \citep[pp.~3763--3764]{bojinov2023}. Asymptotic normality and a conservative variance estimator support testing the weak null and constructing confidence intervals.

The analyst specifies $p$, which need not equal the true order $m$. The main optimal-design analysis sets $p=m$, while Section~4.3 examines misspecification. If $p>m$, inference remains valid under the stated regularity conditions, but the longer required run excludes additional observations and shifts the start of the averaging window. If $p<m$, the estimator is generally biased for the intended persistent-treatment effect. Its asymptotic normality is then centered at its own expectation rather than that causal estimand \citep[p.~3768]{bojinov2023}. The validity of a Fisher test under the sharp no-effect null therefore does not imply valid confidence-interval coverage for the persistent-treatment effect when carryover is underestimated.

\citet{ni2022} is a workshop preprint of \cite{bojinov2023}. It uses 1-carryover assumption.

\cite{liuzhong2026} uses the same m-carryover assumption as \cite{bojinov2023}, and use conditional randomization test with burning early $m$ outputs from each block. 

 

\citet{masoero2023}, a short workshop paper presented at Amazon Science, show why switching can help even when interference is absent. Joint balance over units and periods removes leading contributions from unit and time heterogeneity in their illustrative variance comparison. Their central randomization analysis assumes two potential outcomes per unit--period, with no cross-unit or temporal interference. This paper should therefore be distinguished from work formally accommodating carryover. Its suggestion to space switches beyond a known carryover duration and analyze only subsequent observations is an extension strategy, not the central theorem.

\citet{missault2025} is follow-up paper of \citet{masoero2023}, uses m-carryover assumption (Assumption~4). Assumption~4  requires ``a fixed and given $m$'' in $\{1,\ldots,S-1\}$  They devise \emph{regular balanced switchback designs} which switches between both time and item. Their lagged Horvitz--Thompson analysis requires a lag that covers this order. The paper does design and analysis, including variance estimation and empirical power comparisons. However, its Assumption~1 explicitly excludes cross-item spillovers.  


The adaptive design of \citet{zeng2026} is first developed under no carryover setting and than extended into 1-carryover setting. This is a paper written by Stefan Wager group and airbnb.
  

\citet{chen2026} combine design and analysis by measuring a surrogate $Q_{it}$ that captures the effect of past assignments on the final outcome.  \citep[Assumption~1, p.~4857; PDF p.~4]{chen2026} is $m$-carryover.
 
\subsubsection{Approximate memory requires explicit control of the remaining bias}

\citet{ni2025model} retain nonanticipation while allowing arbitrary-length history. Their exact diminishing-influence restriction is
\begin{equation}\label{eq:ni-history-envelope}
\max_{w_{1:t-q},w'_{1:t-q}}
\left|Y_t(w_{1:t-q},d)-Y_t(w'_{1:t-q},d)\right|
\leq\phi_t(q),\qquad d\in\{0,1\}^{q},
\end{equation}
where $\phi_t(q)$ is decreasing for $q\geq1$ \citep[Assumption~2, Eq.~(8), p.~14]{ni2025model}. This transcribes the source with simplified typography. No particular geometric rate is maintained. Historical modeling instead guides the seasonal cycle, the valid-observation fraction $\rho$, and a temporal-coupling refinement. The design guarantees a predetermined count of balanced usable observations, which directly addresses the instability of random exposure counts.

However, the two-period estimator generally has expectation $\bar\tau=\mathbb{E}(\widehat\tau)$ distinct from the global effect $\tau$. Theorem~5 centers its normal limit at $\mathbb{E}(\widehat\tau)$; the authors explicitly state that the bias remains fixed when this estimator does not adapt to the horizon \citep[p.~25]{ni2025model}. A decreasing influence envelope enables a bias bound, rather than making the resulting interval automatically cover $\tau$.

Variance analysis introduces a further covariance identity,
\begin{equation}\label{eq:ni-covariance-identity}
S_{b,b'}^{1,0}+S_{b,b'}^{0,1}
=S_{b,b'}^{1,1}+S_{b,b'}^{0,0}-S_{b,b'}^{\tau},
\end{equation}
reproduced from Assumption~3 \citep[pp.~18--19]{ni2025model}. The source supplies sufficient conditions, including at most one-period carryover or suitable invariance of treatment contrasts to earlier assignments. These conditions can fail with more complicated historical interactions. Proposition~1 also leaves a conditional-covariance remainder of indeterminate sign in the variance estimate. Fast decay can make it small, but an unspecified decreasing envelope alone does not establish fully conservative standard errors \citep[p.~26]{ni2025model}. Historical autoregressive calibration is therefore useful guidance, while residual global-effect bias and variance approximation remain separate inferential issues.

Finally, \citet{carryover2025} treat an assumed horizon as something to challenge. Their diagnostic theory adopts the following additive finite-lag model, transcribed from Equation~(3) with typography normalized:
\begin{equation}\label{eq:carryover-diagnostic-model}
\begin{aligned}
Y_{i,t}(W_{i,1:t})
 &=\mu_i+\alpha_{i,t}+\delta^{(0)}W_{i,t}\\
 &\quad+\delta^{(1)}W_{i,t-1}+\cdots+\delta^{(m)}W_{i,t-m}+\epsilon_{i,t}.
\end{aligned}
\end{equation}
Here $m$ is the true carryover duration, whereas $\ell$ is the analyst's chosen exclusion lag; their finite-duration assumption excludes effects beyond $m$ \citep[Assumption~2, p.~2; Eq.~(3), p.~3]{carryover2025}. Under fair-coin randomization, differences between aligned estimators with consecutive assumed lags identify the omitted coefficient $-\delta^{(\ell+1)}$ when $\ell<m$. This is an analysis diagnostic, not a new assignment design or a procedure that assumes $m$ is already known. The authors explicitly state that their test cannot confirm a tight upper bound \citep[Section~3.1, p.~3]{carryover2025}. Nonrejection may reflect a weak immediate omitted lag even when more distant effects remain. Across these contributions, credible practice therefore requires recording the chosen physical horizon or decay restriction, distinguishing learned tuning from maintained assumptions, and reporting how conclusions change when that restriction is relaxed.

\subsection{Persistent carryover and models of system dynamics}
\label{sec:dynamic}

A finite carryover window asserts that sufficiently old assignments have exactly zero effect. Dynamic models replace that restriction with assumptions about how information propagates through an evolving system. Their benefits differ: some justify simple estimators without observing the state, whereas others use recorded states or structural equations to improve precision. The associated causal targets also differ. A stationary reward contrast, a finite-horizon global effect, a discounted value contrast, and a local policy derivative should not be interpreted as interchangeable versions of the same estimand.

\subsubsection{Geometric mixing with an unobserved state}

\citet{huwager2025} study design and analysis when the system is a potentially nonstationary Markov decision process. Treatment changes its state transition, and the state can remain unobserved. Their central carryover restriction is Assumption~2 (p.~5): ``There exists a mixing time'' $t_{\mathrm{mix}}<\infty$ such that, for every $t$, $w\in\{0,1\}$, and pair of state distributions $f,f'$,
\begin{equation}
\label{eq:hw-mixing}
 \lVert f'P_t^w-fP_t^w\rVert_{\mathrm{TV}}
 \leq \exp(-1/t_{\mathrm{mix}})\lVert f'-f\rVert_{\mathrm{TV}}.
\end{equation}
This copies their contraction condition; past interventions need never have exactly zero effect. Transition and reward functions may change arbitrarily with time, but the same contraction bound applies throughout. Persistent learning or resource changes that do not dissipate fall outside this restriction.

The distinction between an existence assumption and an operational input is consequential. Assumption~2 does not state that $t_{\mathrm{mix}}$ is observed or known. Nevertheless, their particular rate-guaranteeing prescription $b=(t_{\mathrm{mix}}/2)\log T$ for burn-in length uses its value, or a valid upper bound. State observations and transition estimates are unnecessary for computing the proposed estimators; credible calibration of the memory scale remains necessary for this prescription. Typical trip duration may motivate a choice in transportation applications, but does not by itself certify mixing of all latent behavioral and operational states.

Their December 2025 version also distinguishes the filtered average treatment effect from the full global average treatment effect (GATE). Removing observations immediately after switches changes the periods represented by an ordinary difference-in-means estimator. Under arbitrary temporal effect heterogeneity, this filtered effect can differ from GATE. Near-parametric accuracy for GATE therefore requires their weighted, bias-corrected estimator, or an additional condition making the filtered and global effects sufficiently close. Ordinary burn-in alone is not a general remedy for full-horizon GATE estimation. The paper provides asymptotic inference with conservative variance bounds, as well as design guidance.

\subsubsection{Generalization}

\citet{xiong2024} take a broader design perspective, decomposing error into carryover, periodicity, serially correlated outcomes, and simultaneous interventions. Their Assumption~4.2 requires, for every $t$, a carryover kernel satisfying
\begin{align}
 \int d_t^{\mathrm{co}}(t')f(t')\,dt'&=1,\notag\\
 \delta_t^{\mathrm{co}}(\boldsymbol w_t)
 &=\delta_t^{\mathrm{co}}\int w_{t'}d_t^{\mathrm{co}}(t')f(t')\,dt'.
 \label{eq:xiong-carryover-kernel}
\end{align}
This reproduces the source equations with normalized typography, using $\boldsymbol w_t$ for the assignment history, $f$ for event density, and $\delta_t^{\mathrm{co}}=\delta_t^{\mathrm{co}}(\boldsymbol1_t)$ \citep[Assumption~4.2, p.~21]{xiong2024}. Nonanticipation sets the kernel to zero for $t'>t$. The authors present additive kernel structure as an expository restriction; finite support and geometric tails are examples, rather than one universally imposed window or decay rate. The design is informed by cumulative-effect curves estimated from prior experiments, together with historical outcomes. Their contribution is both design and analysis, principally through bias--variance decomposition and empirical design selection rather than a general new confidence-interval theorem.

This perspective explains why longer intervals need not always improve experiments: they can reduce carryover bias while increasing variance associated with correlated outcomes and treatment allocation. Their empirical Bayes approach makes the role of prior information explicit, but its usefulness for a new policy depends on how well prior response dynamics transfer. The reported nonmonotonic cumulative-effect curves further caution against assuming that all operational effects approach their long-run value smoothly from one side \citep[Sections~1 and~6]{xiong2024}.


\subsubsection{Joint temporal and spatial interference}

\citet{jia2025} extend this approach to clustered switchback experiments, jointly choosing spatial clusters and temporal blocks. Their target is the effect of treating all units throughout rather than controlling all units throughout. Assumption~1 (p.~6) states: ``There exists a constant'' $t_{\mathrm{mix}}>0$ such that, for every individual $i$, time $t$, neighborhood treatment vector $w$, and distributions $f,f'$,
\begin{equation}
 d_{\mathrm{TV}}(fP_{it}^{w},f'P_{it}^{w})
 \leq e^{-1/t_{\mathrm{mix}}}d_{\mathrm{TV}}(f,f').
\end{equation}
The model permits infinite decaying temporal carryover and arbitrary time variation. An interference graph specifies which current assignments affect an individual's reward and transition. Conditional on assignments, individuals' state evolutions are independent; this is more restrictive than unrestricted propagation between neighboring states. Their fractional-neighborhood exposure assumption additionally specifies when a small fraction of discordant neighbors leaves the reward and transition unchanged.

States need not be observed. The graph and exposure probabilities, however, enter the truncated Horvitz--Thompson estimator. Theorem~1 (p.~9) chooses block length and truncation radius $\ell=r=t_{\mathrm{mix}}\log(NT)$, so implementing this particular guarantee requires a defensible mixing bound. Favorable rates depend on graph structure and clustering through exposure probabilities; sparse or spatially structured graphs can yield near-$1/(NT)$ mean squared error, up to logarithmic factors. The main contribution is design and point-estimation error control, rather than a general confidence-interval or hypothesis-testing procedure.

\subsubsection{Observable stationary dynamics and inferential conservatism}

\citet{glynn2020} provide a different route to efficient design and estimation: estimate the unknown stationary transition matrices and select interventions adaptively using the observed state. Their structural assumptions are explicit (p.~2): ``We assume a finite state space'' and ``We assume both'' treatment and control transition matrices ``are irreducible.'' These assumptions impose neither an exact carryover cutoff nor a known numerical mixing time. Irreducibility alone should not be rewritten as the uniform contraction condition in Equation~\eqref{eq:hw-mixing}. Positive limiting state--action sampling frequencies support nonparametric maximum likelihood and asymptotic efficiency.

The knowledge burden shifts to a sufficient observed state, stationarity, and exploration. Unknown transitions and rewards are learned online; they are not assumed available beforehand. The conclusion acknowledges that full state-space knowledge may be infeasible and that asymptotic optimality leaves finite-horizon design questions unresolved. This paper consequently supplies a useful counterexample to the claim that all methods for persistent carryover require a known carryover length or mixing constant.

\citet{blanchet2026} analyze conventional difference-in-means inference under stationary continuous-time dynamics. Assumption~3 (p.~3) specifies, for each policy $\kappa\in\{0,1\}$,
\begin{equation}
 \lVert P_\kappa(x,\cdot;t)-\pi_\kappa\rVert_{\mathrm{TV}}
 \leq C e^{-\lambda_\kappa t},
 \qquad C\geq1,\quad \lambda_\kappa>0.
\end{equation}
This is geometric convergence toward policy-specific stationary distributions, rather than finite memory. Stochastically monotone kernels, stochastic dominance of treatment over control, and monotone rewards are additional substantive restrictions. Under these conditions, the transient effect measured by switching can underestimate the stationary effect. Their Ornstein--Uhlenbeck analysis also shows that conventional Huber--White variance estimates can be too large.

These results diagnose analysis and motivate interval-length and washout choices; they do not establish that conventional inference is always conservative. The geometric constants bound error but are not inputs to the usual difference-in-means calculation. In the illustrative calibration, mean-reversion parameters are estimated. The practical lesson is conditional: structural direction-of-bias results need the monotonicity assumptions, and favorable variance behavior does not remove the estimand discrepancy.

\subsubsection{Reinforcement-learning estimators and the meaning of carryover strength}

\citet{wen2025} compare switch durations for several estimators, including model-based estimation, least-squares temporal difference learning, and double reinforcement learning. They assume independent daily trajectories while allowing time-varying dynamics within each day. Their observed-state Markov condition (p.~4) is
\begin{equation}
 \Pr(S_{t+1}=s'\mid A_t,S_t,\{S_j,A_j,R_j\}_{j<t})
 =p_t(s'\mid A_t,S_t).
\end{equation}
Rewards satisfy conditional mean independence given current state and action, while reward residuals may remain autocorrelated. Their carryover-strength parameter is
\begin{equation}
 \delta=\max_{s,t}\sum_{s'}
 |p_t(s'\mid1,s)-p_t(s'\mid0,s)|.
\end{equation}
It measures treatment-induced changes to transitions, not the last affected lag. In the linear model, Assumption~4 additionally requires $\lVert\Phi_t\rVert_2\leq\rho_\Phi<1$ (p.~22); different estimators use different regularity conditions.

With weak carryover and positively correlated residuals, frequent switching can improve precision by cancelling common temporal fluctuations. Large carryover or predominantly negative residual correlation can favor longer runs. Their workflow calls for assessing state sufficiency, carryover magnitude, and residual dependence. Coefficients, values, and weighting functions are estimated; $\delta$ and stability constants are theoretical conditions rather than asserted known design inputs. Confidence intervals receive empirical evaluation, while the principal contribution is comparative estimation and design efficiency. Neither the preferred cadence nor the requisite observed-state assumption is universal.

\citet{farias2022} instead analyze a given Bernoulli experiment. Their Differences-in-Q estimator targets a stationary effect using observed state--action--reward data and estimated on-policy values. Assumption~1 (pp.~11--12) states
\begin{equation}
 d_{\mathrm{TV}}(P_{1/2}^{k}(s,\cdot),\rho_{1/2})
 \leq C\lambda^k,
\end{equation}
where the experiment's mixed-policy chain has geometric decay. A separate bound $d_{\mathrm{TV}}(p(s,1,\cdot),p(s,0,\cdot))\leq\delta$ controls treatment perturbations. The estimator trades second-order bias for potentially large variance reductions. Theorem~2 (p.~12) centers its central limit theorem at the stationary-experiment expectation of the DQ estimator, rather than the true sustained-policy effect. This approximation bias therefore need not vanish for a fixed intervention. The mixing constants govern guarantees, while learning the value function and observing a sufficient state are the practical requirements.

\citet{shi2023} develop analysis and sequential hypothesis testing for a broader range of assignment rules. Their stationary Markov and conditional-mean assumptions (p.~10), written with normalized notation, are
\begin{align}
 \Pr(S_{t+1}\in B\mid A_t,S_t,H_{t-1})&=P(B;A_t,S_t),\\
 \mathbb{E}(Y_t\mid A_t,S_t,H_{t-1})&=r(A_t,S_t).
\end{align}
Here the observed state must contain important mediators of past treatment effects. Condition~C1(ii) (p.~46) requires behavior-policy geometric ergodicity, written in normalized notation as
\begin{equation}
 \lVert\Pr(S_t\in\cdot\mid S_0=s)-\Pi(\cdot)\rVert_{\mathrm{TV}}
 \leq M(s)\rho^t,\qquad 0\leq\rho<1,
\end{equation}
with $\int M(s)\Pi(ds)<\infty$; alternating designs use analogous even/odd subsequence conditions. These are existence restrictions, not stated known numerical inputs. Their estimand integrates a difference of discounted value functions over a reference initial-state distribution. The analyst's discount factor weights future utility; it is not an assumed physical decay rate for carryover. Thus their sequential tests concern a different causal target from finite-horizon GATE.

\subsubsection{Structural models under congestion and partial observability}

\citet{li2026} study both analysis and design in a queue with state-dependent arrivals. Their main estimand is the derivative of stationary throughput with respect to price, a local total policy effect. Carryover operates through queue occupancy, without an exact cutoff. Assumption~1 (p.~6) requires a bounded queue capacity $K=\min\{k:\lambda_k(p)=0\}$. For interval switchbacks, Assumption~4 (p.~10) requires
\begin{equation}
 \ell_T/T\longrightarrow0,
 \qquad \ell_T/\sqrt{T}\longrightarrow\infty.
\end{equation}
Increasing run lengths control boundary effects; a known numerical mixing time is not stipulated. Small price perturbations, smooth arrival rates, and structural queue assumptions support their central limit theorems and confidence intervals. State-specific direct-effect estimators use recorded queue occupancies and arrivals, and can also work with user-level randomization. Efficiency gains therefore arise from structural information, not a general solution for arbitrary carryover. Their nonstationary extension imposes regime structure.

\citet{sun2024} address partial observability using a controlled ARMA model (p.~12):
\begin{equation}
 Y_t=\mu+\sum_{j=1}^{p}a_jY_{t-j}+bU_t
       +\sum_{j=0}^{q}\theta_j\epsilon_{t-j}.
\end{equation}
Assumption~1 (pp.~14--15) requires all roots of $1-\sum_{j=1}^{p}a_j y^j$ to lie outside the unit circle, and the innovations are independent and identically distributed with zero mean and finite variance. Stable autoregression permits infinite treatment propagation: finite model order $p$ is not a finite carryover horizon. Moving-average dependence represents partial observability. Their design and estimation theory uses a small-signal asymptotic approximation. Crucially, orders and coefficients are not assumed known: historical estimation, AIC/BIC selection, or sequential updating are proposed. The remaining burden is a credible stable linear approximation; the discussion acknowledges that design optimality need not hold outside the small-signal regime and leaves spatial extensions for future work.

\subsubsection{Learning assignments and randomizing treatment durations}

\citet{wu2026} take an estimation procedure as input and optimize its mean squared error through transformer reinforcement learning. The allocation can use the entire observed history, without a specified finite carryover horizon or current-state Markov restriction. Their impossibility result exhibits processes for which restricted-history allocations are suboptimal; it does not assert this for every process. The operative requirement is explicit (p.~7): ``We assume access to a simulation environment'' approximating the experimental data-generating process. Simulated all-treatment and all-control trajectories provide the Monte Carlo target used in training. Learning a simulator from accrued data is proposed when one is unavailable. This design contribution shifts the burden toward counterfactual simulation fidelity and training; it does not establish assumption-free inference or a new general confidence-interval guarantee.

\citet{hu2026} relax the temporal allocation mechanism while retaining finite carryover. Assumption~1 (p.~6) is
\begin{equation}
 w(s)=w'(s)\ \text{for every }s\in[t-m,t]
 \quad\Longrightarrow\quad Y_t(w)=Y_t(w').
\end{equation}
The design alternates treatment labels across random-duration runs. Frequency-domain analysis connects estimator error to temporal imbalance and contamination. A two-rate switching hazard improves the specified worst-case MSE benchmark, but Algorithm~1 explicitly takes $m$ and $T$ as inputs. Randomizing run duration consequently addresses regularity and cadence, not uncertainty about an infinite carryover tail.

\subsubsection{Short reward trajectories and residual bias}

\citet{johari2025} analyze a Bernoulli experiment using truncated policy-gradient estimation. Their Assumption~1 (p.~5) copies a uniform contraction requirement:
\begin{equation}
 \lVert f'P_z^t-fP_z^t\rVert_{\mathrm{TV}}
 \leq\gamma\lVert f'-f\rVert_{\mathrm{TV}},\qquad0<\gamma<1.
\end{equation}
The truncation length $k$ is an analysis choice, not a true last affected lag. The estimator uses observed rewards and assignments without reconstructing the state. Its error includes both geometrically decreasing truncation bias and policy-gradient approximation bias controlled by treatment-induced kernel differences. Larger $k$ does not generally eliminate the latter.

Their inferential result requires more than arbitrary nonstationarity. Assumption~3 (p.~13) imposes the mean-stability condition
\begin{equation}
 T^{-1}\sum_{t=1}^{T}(\mu_t-\bar\mu_T)^2=o(T^{-2/3}),
\end{equation}
where $\mu_t$ is the mean of the estimator's reward-weighted summand and $\bar\mu_T$ its time average, along with a positive limiting variance. Theorem~2 centers the fixed-$k$ central limit theorem at $\mathbb{E}(\widehat\tau_k)$, not GATE. Their consistent HAC variance estimator therefore quantifies sampling variability, while coverage of GATE additionally requires negligible bias or a justified bias allowance. The proposed data-driven truncation heuristic is useful calibration guidance; it does not certify an unknown carryover bound.


\section{Inferential targets, diagnostics, and practical precision}\label{sec:inference}

\subsection{Dynamic effects need not be sustained-policy effects}

An experiment can identify a useful dynamic contrast without identifying the effect of permanently treating the system. \citet{bojinovshephard2019} define time-averaged, lagged causal effects that condition on realized earlier history and average over selected subsequent assignments. Their main framework allows arbitrary dependence on past treatments. The panel generalization of \citet{bojinovpanel2021} retains this flexibility and averages across units and time. A chosen contrast lag therefore need not be an asserted upper bound on causal persistence. These results should not be interpreted as unrestricted identification of the always-treatment versus always-control contrast from a short experiment.

Two distinct forms of non-anticipation matter. An outcome restriction says future treatments cannot affect current potential outcomes. A randomization restriction says assignment cannot depend on unobserved or future potential outcomes. In the local version of \citet[Assumption~2, p.~6]{bojinovshephard2019}, the latter is
\begin{align}
 &\Pr(W_t=w_t\mid W_{1:t-1}=w_{1:t-1},Y_{1:T}(\bullet))\notag\\
 &\qquad=\Pr(W_t=w_t\mid W_{1:t-1}=w_{1:t-1},
 Y_{1:t-1}(w_{1:t-1})).\label{eq:assignment-nonanticipation}
\end{align}
This permits assignment based on observed history. It is different from requiring independently randomized treatments. Their outcome notation already restricts potential outcomes to current and past assignments. Their optional ``$m$-period causal impact'' restriction appears separately in Appendix~B.2; it is not necessary for every estimand or inferential result in that framework.

The panel model indexes outcomes by each unit's own treatment path and thus excludes interference between units, while permitting arbitrary within-unit temporal dependence \citep[Definition~1 and Assumption~1, p.~5]{bojinovpanel2021}. It provides exact randomization tests for sharp nulls and conservative asymptotic inference for average effects. It also shows why ordinary fixed-effects regressions can fail to recover contemporaneous causal effects when carryover combines with persistent assignment. Neither repeated measurement nor unit fixed effects alone resolve this problem.

\subsection{Partially sharp nulls and what diagnostics establish}

\citet{liuzhong2026} study inference for a total-effect null rather than assuming every treatment history has the same outcome. Their temporal restrictions are explicitly separate. Reproducing their mathematical statements, Assumption~2.1 requires
\[
 Y_t(w)=Y_t(w')\quad\text{whenever }w_{1:t}=w'_{1:t},
\]
whereas Assumption~2.2 fixes $m\in\{0,\ldots,T-1\}$ and requires, for $t\geq m+1$,
\[
 Y_t(w)=Y_t(w')\quad\text{whenever }w_{t-m:T}=w'_{t-m:T}.
\]
Both statements quantify over all treatment paths \citep[p.~6]{liuzhong2026}. The second holds future assignments fixed and therefore does not itself exclude anticipation. Together, they imply dependence only on $w_{t-m:t}$.

With $Y_t^{[m]}$ denoting the resulting local potential outcome, their null is
\[
 H_0^{\mathrm{tot}}(m):\quad
 Y_t^{[m]}(1_{m+1})=Y_t^{[m]}(0_{m+1}),
 \qquad t=m+1,\ldots,T.
\]
This is a period-specific, \emph{partially sharp} null: it equates the two constant histories but leaves outcomes under mixed histories unrestricted. Consequently, holding observed outcomes fixed while rerandomizing the entire assignment path is generally unjustified. Their conditional randomization test instead pools predetermined blocks into sections, retains sections that happened to be entirely treated or entirely controlled, and excludes the first $m$ periods. The null then imputes the counterfactuals needed for the restricted randomization distribution \citep[Sections~2--3]{liuzhong2026}.

Finite-sample validity requires a correctly specified, prespecified upper bound $m$. Longer exclusions and larger sections can protect against carryover while reducing the informative randomized comparisons and worsening attainable $p$-values. Tests of average-effect weak nulls allow effects of opposite signs to cancel; they do not impute missing outcomes. The paper instead establishes asymptotic validity through studentization under additional regularity, equal-length blocks, $m<L$, and an analysis without assignment-dependent selection.

The carryover diagnostic tests proposed horizons while maintaining outcome non-anticipation. Sequentially testing nested horizons from shortest to longest controls erroneous rejection of valid horizons. Its inferential direction is crucial: rejected horizons supply a lower confidence bound on persistence. The source states, ``The first nonrejected value determines where the procedure stops; it should not be interpreted as an estimate of the carryover horizon'' \citep[Section~4.1, p.~14]{liuzhong2026}. Nonrejection does not certify a sufficiently long washout, and the paper does not justify selecting a horizon and testing the total effect with the same outcomes. Anticipation receives a separate pairwise randomization diagnostic. These are analysis tools with design implications, rather than a replacement for scientific knowledge of plausible persistence.

\subsection{Inference with decaying rather than truncated effects}

\citet{linding2025} justify regression-based estimation of lagged effects and conservative heteroskedasticity-and-autocorrelation-consistent (HAC) standard errors under a misspecified working regression. Their main assignment assumption is independent $Z_t$ with known probabilities $p_t$, bounded away from zero and one. Their outcome no-anticipation assumption still allows dependence on all previous assignments.

Their Assumption~3.2 defines the following expected absolute coupling discrepancy, reproduced in source notation:
\begin{align}
 \theta_{T,k}=\max_{k<t\leq T}\mathbb{E}\big[\big|
 &Y_t(Z_{t:t-k+1},Z_{t-k:1})\notag\\[-2pt]
 &-Y_t(Z_{t:t-k+1},Z'_{t-k:1})\big|\big],
\end{align}
where $Z'_{t-k:1}$ is an independent copy of the older assignment history. Its two conditions are
\[
 T^{-1}\sum_{k=1}^T\theta_{T,k}\longrightarrow0,
 \qquad
 \theta_{T,k}=O(k^{-1-\delta})\quad\text{for some }\delta>0.
\]
The weaker condition supports consistency, with additional restrictions if the regression dimension grows; the polynomial decay condition supports asymptotic normality \citep[Assumption~3.2, pp.~8--9]{linding2025}. This is a causal-history coupling restriction, not a claim that observed autocorrelation alone measures carryover.

An alternative is their Assumption~3.3:
\[
 Y_t(Z_{t:1})=Y_t(Z_{t:t-m}),\qquad m<t\leq T,
 \qquad m/T\to0,\quad m^4/T\to0.
\]
Here $m=m(T)$ can grow. The two rates serve different consistency and normality results. The displayed window contains current treatment and $m$ prior assignments. Regression lag count $K$ is a selected target dimension, not necessarily the true cutoff. A numerical decay exponent is not an input to ordinary least squares, but the decay property remains a maintained assumption. HAC bandwidth $L$ must grow suitably: under their finite-memory regime the variance result includes $Km^2/L\to0$ and $K^4L^2/T\to0$ \citep[Theorem~3.3, p.~13]{linding2025}. Unknown dependence can therefore make bandwidth choice consequential even when $m$ is not supplied directly to the estimator.

These lagged effects are independent of the realized history, but generally depend on the assignment probabilities used to average over other histories. Design invariance holds in special additive models \citep[Remark~2.1, p.~8]{linding2025}. Working-model robustness should not be expanded into robustness to arbitrary assignment dependence or arbitrarily persistent causal effects.

\citet{liangrecht2025} instead impose linear time-invariant dynamics. Their structural model is
\[
 y_t=\sum_{s=0}^t x_sg_{t-s}+e_t,
\]
with independently randomized binary inputs and a baseline sequence $e_t$ oblivious to randomization. Additivity and an invariant impulse response $g$ replace unrestricted history dependence. Persistent effects are permitted; geometric decay is a sufficient regularity condition rather than an exact finite cutoff. Their actual Assumption~1 concerns both $g$ and the target weights $q$. Writing $\circledast$ for the source's circular convolution notation, it requires
\begin{align}
 &\big\lVert(|g|\circledast|q|)\circledast(|g|\circledast|q|)\big\rVert_2^2\notag\\
 &\quad\leq CT^{1-\epsilon}\Big(\lVert g\circledast q\rVert_2^2
   +\langle g\circledast g,q\circledast q\rangle
   -2\langle q,g\rangle^2\Big)^2,
 \label{eq:liangrecht-normality}
\end{align}
for universal constants $C>1$ and $\epsilon>0$ \citep[Assumption~1, p.~15]{liangrecht2025}. This transcribes the source inequality with normalized convolution typography. It controls the quadratic term's asymptotic normality; summability alone does not replace this condition. Exponential decay $|g_t|\lesssim\rho^t$, $0<\rho<1$, together with nondegenerate variance and suitably limited target support, is a sufficient route to this condition. Separately, the uniform response-estimation result assumes $\lVert g\rVert_1\leq C$ and $|e_t|\leq M$ for fixed constants \citep[Theorem~3, p.~17]{liangrecht2025}; those bounds are not Assumption~1.

The decay coefficient need not be numerically known to form the moment estimator. Nevertheless, choosing a truncated cumulative target and controlling its omitted tail are different tasks. The presented confidence-interval construction uses targets supported on $K=O(\log T)$ entries. Unknown slow decay threatens approximation to a full policy effect. Moreover, a finite-horizon always-treatment contrast includes boundary weights $(T-t)/T$ on $g_t$; it is not automatically an unweighted truncated sum. Linear-response structure, stable effects, and target complexity remain substantive limitations.

This distinction also clarifies what ``known persistence'' means. For instance, an assumed envelope $|g_t|\leq C\rho^t$ implies the elementary bound $\sum_{t>K}|g_t|\leq C\rho^{K+1}/(1-\rho)$. This implication is a review calculation, rather than an additional result claimed by the cited paper. Knowing only that some $\rho<1$ exists does not numerically bound a finite experiment's omitted tail. A quantitative approximation guarantee needs a defensible envelope or a sensitivity analysis over plausible envelopes. Neither a visually small sample autocorrelation nor a nonsignificant lag coefficient supplies that counterfactual bound.

\subsection{Model order and sequence support determine feasible inference}

\citet{guoliang2026} emphasize design optimization for cumulative or lag-specific effects under a finite-order impulse-response model. Their source specification is
\[
 Y_t=\alpha+\sum_{k=1}^K\beta_kZ_{t-k+1}+\epsilon_t,
 \qquad \epsilon_t\stackrel{\mathrm{iid}}{\sim}N(0,\sigma^2),
\]
with errors independent of assignments and $K$ fixed as $T$ grows \citep[Equation~(3.1), pp.~6--7]{guoliang2026}. Equivalently,
\[
 Y_t(z_1,\ldots,z_T)=Y_t(z_{t-K+1},\ldots,z_t).
\]
Thus $K$ includes current treatment and allows $K-1$ lagged periods. Design optimization must specify this order; the paper does not provide automatic recovery of an unknown causal cutoff. Their Markov assignment framework compares random-switch and cycle-switch classes. Preferred persistence depends on the target, so optimality is conditional on the model and candidate class. Simulated misspecification performance does not remove the finite-order, additive structure from the optimality results.

Crossover analysis makes the identification role of sequence support especially clear. \citet[Assumption~2, p.~6]{jiangding2025} impose, when scientifically appropriate,
\[
 Y_{it}(\vec z)=Y_{it}(\vec z')
 \quad\text{if}\quad
 \vec z_{[t-k+1,t]}=\vec z'_{[t-k+1,t]}.
\]
Their $k=1$ means no lagged carryover; it corresponds to $m=0$ in the preceding finite-memory notation. No anticipation is separate, and time-invariance of treatment-sequence effects is another restriction. With all four two-period sequences, instantaneous and carryover contrasts can be identified without imposing no carryover; with only $AB$ and $BA$, unrestricted second-period direct and carryover effects are aliased. The assumed order can therefore determine identification, rather than merely improve precision. Restricted weighted least squares and conservative design-based variance estimates tolerate working-model misspecification, while the encoded causal restrictions remain essential.

The source motivates finite memory through short-acting interventions and sufficiently spaced periods, while warning that learning, delayed responses, and cumulative biological or behavioral effects can persist. These are reasons to justify restrictions from the intervention and observation clock, rather than from regression convenience. Likewise, concealed assignments can support no anticipation, whereas predictable or announced future prices can undermine it. Restricting an outcome model for efficient estimation and excluding possible causal histories for identification are materially different commitments.

The broader habituation designs of \citet{basse2023} also show why finite memory need not be universal. Their main pulse/wedge allocation designs allow indefinite persistence and exclude interference across units. Only the optional recycling extension requires
\[
 Y_{it}(e_{t'})=Y_{it}(0)\quad\text{for }0<t'\leq t-k,
\]
where $e_{t'}$ is a single treatment pulse \citep[Assumption~3, p.~15]{basse2023}. This is a pulse-specific restriction, rather than equality for every pair of histories sharing a recent window. A justified $k$ permits recycling previously exposed units as controls; uncertainty about $k$ threatens that efficiency gain without invalidating every baseline habituation contrast.

\subsection{Spatial exposure restrictions and practical precision}

\citet{nipanel2024} explicitly handle both spatial and temporal interference through balanced temporal assignment and randomized spatial clustering. Their local working paper restricts spillovers to a known neighborhood $\mathcal N(i)$ on a two-dimensional lattice. Using source notation with $\mathcal N(i)$ for the neighborhood, Assumptions~1--2 require
\begin{align*}
 Y_{it}(w)&=Y_{it}(w')
 &&\text{if }w_{\mathcal N(i),1:T}=w'_{\mathcal N(i),1:T},\\
 Y_{it}(w)&=Y_{it}(w')
 &&\text{if }w_{1:N,t-1:t}=w'_{1:N,t-1:t}.
\end{align*}
The second equality restricts dependence to current and previous periods and also excludes anticipation \citep[p.~8]{nipanel2024}. Choosing a period long enough to encompass physical persistence relies on an external upper bound; changing the clock does not prove that a tail vanishes. Domain knowledge must likewise justify spatial neighborhood coverage. Their design and confidence-interval guarantees concern these exposure restrictions; broader-network simulations provide sensitivity evidence rather than replacing them.

Practical precision studies address another question. \citet{ferrari2026} compare daily and weekly switching in airline ancillary pricing using a common contemporaneous-effect fixed-effects model. They explicitly omit a causal carryover mechanism from their generator and identify this as a limitation. Their calibrated serial correlation is outcome dependence, not proof of a washout horizon. Reported precision gains are consequently context specific.

\citet{pankratev2026power} derive an asymptotic power approximation under a contemporaneous treatment model with cluster, time, cluster-time, and individual noise components. Cluster-size imbalance amplifies macro-shock variance, so adding observations within existing cells cannot substitute for independent assignment information. The model does not establish inference robust to treatment-induced persistence or cross-cluster spillovers.

\citet{pankratev2026variance} compare adjustment methods by simulation, including both carryover and spatial contamination. Their carryover generator is explicitly
\[
 C_{jh}=\rho_C\tau_j\sum_{k=1}^{3}w_k(T_{j,h-k}-T_{jh}),
 \qquad w=(0.3,0.2,0.1),
\]
rather than an identified three-period persistence bound in field data \citep[Equation~(3), p.~6]{pankratev2026variance}. Adjustment improves precision in favorable regimes but cannot repair interference bias; sufficiently precise intervals can support wrong-sign decisions when contamination dominates. Results depend on the selected generator, covariate information, and cluster count. These studies motivate power calibration and sensitivity analysis while reinforcing the distinction between estimating an effect precisely and identifying the intended effect.


\section{Rollout decisions under temporal restrictions}\label{sec:decisions}

The inferential endpoint can itself change the appropriate assumptions and objective. \citet{ni2025decision} take an experimental design as given and compare worst-case simple regret from adopting treatment or retaining control. Their baseline bounded-outcome uncertainty set yields a familiar inverse-propensity comparison. Refinement uses the temporal restriction
\begin{equation}\label{eq:decision-smoothness}
 |Y_{t+1}(w)-Y_t(w)|\leq L,
\end{equation}
reproduced from Assumption~3 \citep[Eq.~(10), p.~10]{ni2025decision}. This constrains adjacent outcomes under the same assignment, rather than saying that the treatment effect itself is constant. Historical data can inform $L$ if treatment does not materially change smoothness; that transfer is a scientific assumption, not a consequence of randomization.

Their carryover extension chooses a neighborhood $b$ and imposes diminishing effects of a historical treatment perturbation. One of the two inequalities in Assumption~4 is
\begin{align}\label{eq:decision-monotonicity}
 &\left|Y_t(w,\mathbf{1}_{m},\mathbf{0}_{b+1-m})
       -Y_t(w,\mathbf{1}_{m-1},\mathbf{0}_{b+2-m})\right|\notag\\
 &\quad\geq
 \left|Y_{t+1}(w,\mathbf{1}_{m},\mathbf{0}_{b+2-m})
       -Y_{t+1}(w,\mathbf{1}_{m-1},\mathbf{0}_{b+3-m})\right|,
 \qquad 1\leq m\leq b.
\end{align}
This copies Equation~(31) with typography normalized; the second inequality swaps zeros and ones throughout \citep[p.~21]{ni2025decision}. Here $m$ indexes perturbations within the chosen neighborhood: it is not an assumed known maximum lag. The inequalities require declining absolute influence but specify neither a geometric rate nor an exact last affected period.

These restrictions allow transition observations to constrain unobserved outcomes in a robust optimization problem. They should still be assessed against the intervention: abrupt peak demand can challenge smoothness, while delayed feedback can challenge monotonic perturbation effects. A selected neighborhood also approximates permanent exposure. The paper's simulation comparisons calibrate a risk-preference parameter to match false-positive rates in the scenarios studied; they do not establish an unrestricted finite-sample testing guarantee. The contribution is a decision rule under an uncertainty set, with conventional inference serving a complementary role.

\section{Critical synthesis and research priorities}\label{sec:synthesis}

\subsection{What normality and precision leave unresolved}

A useful common decomposition, introduced here as review notation, is
\begin{equation}\label{eq:error-decomposition}
 \widehat\tau-\tau_T
 =\underbrace{\widehat\tau-\mathbb{E}(\widehat\tau)}_{\text{assignment or sampling fluctuation}}
 +\underbrace{\mathbb{E}(\widehat\tau)-\tau_T}_{\text{target discrepancy}}.
\end{equation}
A variance estimator and central limit theorem address the first term. Exact memory, correctly specified exposure, controlled truncation tails, or a justified structural model address the second. When effects remain after the chosen window, increasing precision can sharpen an estimate around the wrong target. Model-assisted balancing, policy-gradient approximations, and practical adjustment studies all make this distinction visible in different ways.

For an interval intended to cover $\tau_T$, the discrepancy must vanish relative to the standard error, or be covered by a defensible bias allowance. If a valid fluctuation interval has half-width $a_T$ and a scientifically justified deterministic bound is $|\mathbb{E}(\widehat\tau)-\tau_T|\leq B_T$, enlarging its half-width to $a_T+B_T$ is a direct conservative implication. This elementary review calculation does not estimate $B_T$ from the data. If the bound is itself uncertain, that uncertainty also requires treatment. Reporting standard errors alone cannot substitute for it.

\subsection{Which assumptions are credible, and for which intervention?}

No single carryover model is uniformly most realistic. Exact memory can be credible for a known processing deadline but implausible for durable learning. Geometric mixing permits persistent tails, yet uniform contraction across states, policies and times can be demanding. An observed-state method avoids specifying a last lag only if the measured state captures the pathways carrying earlier effects. Linear impulse-response models permit long histories while ruling out nonlinear interactions and time-varying response kernels. Simulator-based optimization shifts attention toward whether a fitted simulator predicts counterfactual sustained-policy trajectories.

These are alternative information commitments, not a hierarchy in which weaker-looking notation always implies safer inference. Historical seasonality can support balancing without revealing how a novel policy changes memory. A simulator reproducing baseline autocorrelation can still miss policy-induced behavioral adaptation. Calibration should therefore separate untreated outcome dependence, treatment-induced transition changes, and the longer-run response relevant to deployment.

Spatial aggregation creates a related tradeoff. Larger clusters can better contain interference but leave fewer independently randomized units and greater heterogeneity. Smaller clusters increase nominal replication while requiring stronger exposure modeling. The known neighborhood in a finite spatial model and the graph-based exposure probabilities in a clustered mixing model require distinct justifications. Cluster-level standard errors alone do not remove contamination across cluster boundaries.

\subsection{A reporting sequence and open problems}

A defensible report should state the intervention, unit, observation clock, initial condition, and causal target; separate assignment from the carryover or exposure model; and identify required bounds, estimated nuisance quantities, and their data sources. It should distinguish estimates, intervals, sharp and weak null tests, diagnostics, and deployment decisions. Sensitivity to longer memory, alternative decay envelopes, and spatial contamination should accompany precision claims. This reporting sequence is our synthesis.

Several research needs follow. Horizon diagnostics need power guarantees against distant, weak, nonlinear, or cancelling effects and valid inference after selection. Model-assisted designs need bias-aware intervals under unknown persistent carryover. Adaptive designs need methods combining long memory and cross-unit exposure. Structural and simulator-based methods need checks of intervention-induced dynamics beyond baseline fit. Comparisons should align causal targets, initial conditions, observed state information, and pilot-data and computational costs before recommending a cadence or estimator.

\section{Conclusion}

Switchback research now spans fixed-window designs, persistent dynamics, adaptive allocation, diagnostics, and rollout decisions. Its central practical challenge is connecting short experimental exposure to the policy effect of interest. Finite memory, decay restrictions, observed states, and simulators supply that connection under different information commitments. A credible analysis should state those commitments, distinguish sampling uncertainty from target discrepancy, and report sensitivity to temporal and spatial assumptions. Better precision alone cannot establish that an experiment identifies the effect of sustained deployment.

\clearpage
\begin{singlespace}
\small
\bibliography{switchback-review}
\end{singlespace}
\end{document}
